% Gemini theme
% https://github.com/anishathalye/gemini

\documentclass[final]{beamer}

% ====================
% Packages
% ====================

\usepackage{geometry}
\geometry{
  a4paper,
  total={170mm,257mm},
  left=20mm,
  top=20mm,
}
\usepackage{setspace}
\usepackage{graphicx}
\usepackage{xcolor}
\usepackage{kotex}
\usepackage{csquotes}
\usepackage{tabularray}
\usepackage{relsize}

%%% Typesetting for listings
\usepackage{listings}

\definecolor{dkgreen}{rgb}{0,0.6,0}
\definecolor{ltblue}{rgb}{0,0.4,0.4}
\definecolor{dkviolet}{rgb}{0.3,0,0.5}

%%% lstlisting coq style (inspired from a file of Assia Mahboubi)
\lstdefinelanguage{Coq}{ 
    % Anything betweeen $ becomes LaTeX math mode
    mathescape=true,
    % Comments may or not include Latex commands
    texcl=false, 
    % Vernacular commands
    morekeywords=[1]{Section, Module, End, Require, Import, Export,
        Variable, Variables, Parameter, Parameters, Axiom, Hypothesis,
        Hypotheses, Notation, Local, Tactic, Reserved, Scope, Open, Close,
        Bind, Delimit, Definition, Let, Ltac, Fixpoint, CoFixpoint, Add,
        Morphism, Relation, Implicit, Arguments, Unset, Contextual,
        Strict, Prenex, Implicits, Inductive, CoInductive, Record,
        Structure, Canonical, Coercion, Context, Class, Global, Instance,
        Program, Infix, Theorem, Lemma, Corollary, Proposition, Fact,
        Remark, Example, Proof, Goal, Save, Qed, Defined, Hint, Resolve,
        Rewrite, View, Search, Show, Print, Printing, All, Eval, Check,
        Projections, inside, outside, Def},
    % Gallina
    morekeywords=[2]{forall, exists, exists2, fun, fix, cofix, struct,
        match, with, end, as, in, return, let, if, is, then, else, for, of,
        nosimpl, when},
    % Sorts
    morekeywords=[3]{Type, Prop, Set, true, false, option},
    % Various tactics, some are std Coq subsumed by ssr, for the manual purpose
    morekeywords=[4]{pose, set, move, case, elim, apply, clear, hnf,
        intro, intros, generalize, rename, pattern, after, destruct,
        induction, using, refine, inversion, injection, rewrite, congr,
        unlock, compute, ring, field, fourier, replace, fold, unfold,
        change, cutrewrite, simpl, have, suff, wlog, suffices, without,
        loss, nat_norm, assert, cut, trivial, revert, bool_congr, nat_congr,
        symmetry, transitivity, auto, split, left, right, autorewrite},
    % Terminators
    morekeywords=[5]{by, done, exact, reflexivity, tauto, romega, omega,
        assumption, solve, contradiction, discriminate},
    % Control
    morekeywords=[6]{do, last, first, try, idtac, repeat},
    % Comments delimiters, we do turn this off for the manual
    morecomment=[s]{(*}{*)},
    % Spaces are not displayed as a special character
    showstringspaces=false,
    % String delimiters
    morestring=[b]",
    morestring=[d],
    % Size of tabulations
    tabsize=3,
    % Enables ASCII chars 128 to 255
    extendedchars=false,
    % Case sensitivity
    sensitive=true,
    % Automatic breaking of long lines
    breaklines=false,
    % Default style fors listings
    basicstyle=\small,
    % Position of captions is bottom
    captionpos=b,
    % flexible columns
    columns=[l]flexible,
    % Style for (listings') identifiers
    identifierstyle={\ttfamily\color{black}},
    % Style for declaration keywords
    keywordstyle=[1]{\ttfamily\color{dkviolet}},
    % Style for gallina keywords
    keywordstyle=[2]{\ttfamily\color{dkgreen}},
    % Style for sorts keywords
    keywordstyle=[3]{\ttfamily\color{ltblue}},
    % Style for tactics keywords
    keywordstyle=[4]{\ttfamily\color{dkblue}},
    % Style for terminators keywords
    keywordstyle=[5]{\ttfamily\color{dkred}},
    %Style for iterators
    %keywordstyle=[6]{\ttfamily\color{dkpink}},
    % Style for strings
    stringstyle=\ttfamily,
    % Style for comments
    commentstyle={\ttfamily\color{dkgreen}},
    %moredelim=**[is][\ttfamily\color{red}]{/&}{&/},
    literate=
    {\\forall}{{\color{dkgreen}{$\forall\;$}}}1
    {\\exists}{{$\exists\;$}}1
    {\ge -}{{$\geftarrow\;$}}1
    {=>}{{$\Rightarrow\;$}}1
    {==}{{\code{==}\;}}1
    {==>}{{\code{==>}\;}}1
    %    {:>}{{\code{:>}\;}}1
    {->}{{$\rightarrow\;$}}1
    {\ge ->}{{$\geftrightarrow\;$}}1
    {\ge ==}{{$\geq\;$}}1
    {\#}{{$^\star$}}1 
    {\\o}{{$\circ\;$}}1 
    {\@}{{$\cdot$}}1 
    {\/\\}{{$\text{ and }\;$}}1
    {\\\/}{{$\vee\;$}}1
    {++}{{\code{++}}}1
    {~}{{\ }}1
    {\@\@}{{$@$}}1
    {\\mapsto}{{$\mapsto\;$}}1
    {\\hline}{{\rule{\linewidth}{0.5pt}}}1
    %
}[keywords,comments,strings]

%%% Math settings
\usepackage{amssymb,amsmath,amsthm,mathtools}
\usepackage[math-style=TeX,bold-style=TeX]{unicode-math}
\theoremstyle{definition}
\newtheorem{definition}{Definition}[section]
\newtheorem{example}{Example}[section]
\newtheorem{lem}{Lemma}[section]
\newtheorem{thm}{Theorem}[section]
\newtheorem{cor}{Corollary}[section]
\newtheorem{clm}{Claim}[section]

%%% Font settings
%\setmainfont{Libertinus Serif}
%\setsansfont{Libertinus Sans}[Scale=MatchUppercase]
%\setmonofont{JuliaMono}[Scale=MatchLowercase]
%\setmainhangulfont{Noto Serif CJK KR}
%\setmonohangulfont{D2Coding}

%%% PL constructs
\usepackage{galois}
\usepackage{ebproof}
%% InfRule command
\ExplSyntaxOn
\NewTblrEnviron{@ruleenv}
\SetTblrInner[@ruleenv]{belowsep=0pt,stretch=0}
\SetTblrOuter[@ruleenv]{baseline=b}
\NewDocumentEnvironment { InfRule } { m +b }
  {
    \begin{@ruleenv}{l}
      \smaller{\textsc{#1}} \\
      \begin{prooftree} #2 \end{prooftree}
    \end{@ruleenv}
  }
  {}
\ExplSyntaxOff

%% Modification for \infer
\ExplSyntaxOn
\int_new:N \g__ebproof_sublevel_int
\box_new:N \g__ebproof_substack_box
\seq_new:N \g__ebproof_substack_seq

\cs_new:Nn \__ebproof_clear_substack:
  {
    \int_gset:Nn \g__ebproof_sublevel_int { 0 }
    \hbox_gset:Nn \g__ebproof_substack_box { }
    \seq_gclear:N \g__ebproof_substack_seq
  }

\cs_new:Nn \__ebproof_subpush:N
  {
    \int_gincr:N \g__ebproof_sublevel_int
    \hbox_gset:Nn \g__ebproof_substack_box
      { \hbox_unpack:N \g__ebproof_substack_box \box_use:c { \__ebproof_box:N #1 } }
    \seq_gput_left:Nv \g__ebproof_substack_seq
      { \__ebproof_marks:N #1 }
  }

\cs_new:Nn \__ebproof_subpop:N
  {
    \int_compare:nNnTF { \g__ebproof_sublevel_int } > { 0 }
      {
        \int_gdecr:N \g__ebproof_sublevel_int
        \hbox_gset:Nn \g__ebproof_substack_box {
          \hbox_unpack:N \g__ebproof_substack_box
          \box_gset_to_last:N \g_tmpa_box
        }
        \box_set_eq_drop:cN { \__ebproof_box:N #1 } \g_tmpa_box
        \seq_gpop_left:NN \g__ebproof_substack_seq \l_tmpa_tl
        \tl_set_eq:cN { \__ebproof_marks:N #1 } \l_tmpa_tl
      }
      { \PackageError{ebproof}{Missing~premiss~in~a~proof~tree}{} \__ebproof_clear:N #1 }
  }

\cs_new:Nn \__ebproof_append_subvertical:NN
  {
    \bool_if:NTF \l__ebproof_updown_bool
      { \__ebproof_append_above:NN #1 #2 }
      { \__ebproof_append_below:NN #1 #2 }
  }

\cs_new:Nn \__ebproof_join_subvertical:n
  {
    \group_begin:
    \__ebproof_subpop:N \l__ebproof_a_box
    \prg_replicate:nn { #1 - 1 }
      {
        \__ebproof_subpop:N \l__ebproof_b_box
        \__ebproof_enlarge_conclusion:NN \l__ebproof_b_box \l__ebproof_a_box

        \__ebproof_make_vertical:Nnnn \l__ebproof_c_box
          { \__ebproof_mark:Nn \l__ebproof_b_box {axis} - \__ebproof_mark:Nn \l__ebproof_b_box {left} }
          { \__ebproof_mark:Nn \l__ebproof_b_box {right} - \__ebproof_mark:Nn \l__ebproof_b_box {left} }
          { \skip_vertical:N \l__ebproof_rule_margin_dim }
        \__ebproof_vcenter:N \l__ebproof_b_box
        \__ebproof_append_subvertical:NN \l__ebproof_a_box \l__ebproof_c_box

        \__ebproof_append_subvertical:NN \l__ebproof_a_box \l__ebproof_b_box
      }
    \__ebproof_push:N \l__ebproof_a_box
    \group_end:
  }

\cs_new:Nn \__ebproof_renew_statement:nnn
  {
    \exp_args:Nc \RenewDocumentCommand { ebproof#1 }{ #2 } { #3 }
    \seq_gput_right:Nn \g__ebproof_statements_seq { #1 }
  }
\__ebproof_renew_statement:nnn { infer } { O{} m O{} m }
  {
    \group_begin:
    \__ebproof_restore_statements:
    \keys_set_known:nnN { ebproof / rule~style } { #1 } \l_tmpa_tl
    \keys_set:nV { ebproof } \l_tmpa_tl
    \tl_set:Nn \l__ebproof_right_label_tl { #3 }

    \__ebproof_clear_substack:
    \clist_map_inline:nn { #2 }
      {
        \__ebproof_join_horizontal:n { ##1 }
        \__ebproof_pop:N \l__ebproof_a_box
        \__ebproof_subpush:N \l__ebproof_a_box
      }
    \__ebproof_join_subvertical:n { \clist_count:n { #2 } }

    \__ebproof_push_statement:n { #4 }
    \__ebproof_join_vertical:
    \group_end:
  }
\ExplSyntaxOff
\ebproofset{left label template=\textsc{[\inserttext]}}
\ebproofset{center=false}

%%% Custom commands
\newcommand*{\vbar}{|}
\newcommand*{\finto}{\xrightarrow{\text{\textrm{fin}}}}
\newcommand*{\istype}{\mathrel{⩴}}
\newcommand*{\ortype}{\mathrel{|}}
\newcommand*{\cons}{::}
\newcommand*{\pset}{\mathscr{P}}
\newcommand*{\closeloc}[2]{\prescript{\text{\textbackslash}{#1}}{}{#2}}
\newcommand*{\openloc}[2]{{#1}^{#2}}

\def\ovbarw{1.2mu}
\def\ovbarh{1}
\newcommand*{\ovbar}[1]{\mkern \ovbarw\overline{\mkern-\ovbarw{\smash{#1}\scalebox{1}[\ovbarh]{\vphantom{i}}}\mkern-\ovbarw}\mkern \ovbarw}
\newcommand*{\A}[1]{\overset{\,_{\mbox{\Large .}}}{#1}}
\newcommand*{\Abs}[1]{{#1}^{\#}}
\newcommand*{\Expr}{\text{Expr}}
\newcommand*{\ExprVar}{\text{Var}}
\newcommand*{\Module}{\text{Module}}
\newcommand*{\ModVar}{\text{ModVar}}
\newcommand*{\modid}{d}
\newcommand*{\Time}{\mathbb{T}}
\newcommand*{\ATime}{\A{\Time}}
\newcommand*{\Ctx}{\text{Env}}
\newcommand*{\ctx}{\sigma}
\newcommand*{\Value}{\text{Val}}
\newcommand*{\Walue}{\text{WVal}}
\newcommand*{\Mem}{\text{Mem}}
\newcommand*{\Loc}{\text{Loc}}
\newcommand*{\FLoc}{\text{FLoc}}
\newcommand*{\Left}{\text{Left}}
\newcommand*{\Right}{\text{Right}}
\newcommand*{\Sig}{\text{Sig}}
\newcommand*{\mem}{m}
\newcommand*{\dom}{\text{dom}}
\newcommand*{\AMem}{\A{\text{Mem}}}
\newcommand*{\State}{\text{State}}
\newcommand*{\AState}{\A{\text{State}}}
\newcommand*{\Result}{\text{Result}}
\newcommand*{\AResult}{\A{\text{Result}}}
\newcommand*{\Tick}{\text{Tick}}
\newcommand*{\lfp}{\mathsf{lfp}}
\newcommand*{\Step}{\mathsf{Step}}
\newcommand*{\semarrow}{\Downarrow}
\newcommand*{\asemarrow}{\A{\rightsquigarrow}}
\newcommand*{\synlink}{\rtimes}
\newcommand*{\semlink}{\mathbin{\rotatebox[origin=c]{180}{$\propto$}}}
\newcommand*{\link}[2]{{#1}\rtimes{#2}}
\newcommand*{\mt}{\mathsf{empty}}
\newcommand*{\valid}{\mathsf{valid}}
\newcommand*{\Path}{\text{Path}}
\newcommand*{\equivalent}{\sim}

\newcommand*{\doubleplus}{\ensuremath{\mathbin{+\mkern-3mu+}}}
\newcommand*{\project}{\text{\texttt{:>} }}
\newcommand*{\Exp}{\mathsf{Exp}}
\newcommand*{\Imp}{\mathsf{Imp}}
\newcommand*{\Fin}{\mathsf{Fin}}
\newcommand*{\Link}{\mathsf{Link}}
\newcommand*{\sembracket}[1]{\lBrack{#1}\rBrack}
\newcommand*{\fin}[2]{{#1}\xrightarrow{\text{fin}}{#2}}
\newcommand*{\addr}{\mathsf{addr}}
\newcommand*{\tick}{\mathsf{tick}}
\newcommand*{\modctx}{\mathsf{ctx}}
\newcommand*{\mapinject}[2]{{#2}[{#1}]}
\newcommand*{\inject}[2]{{#2}\langle{#1}\rangle}
\newcommand*{\deletepre}[2]{{#2}\overline{\doubleplus}{#1}}
\newcommand*{\deletemap}[2]{{#1}\overline{[{#2}]}}
\newcommand*{\delete}[2]{{#2}{\langle{#1}\rangle}^{-1}}
\newcommand*{\filter}{\mathsf{filter}}
\newcommand*{\Lete}{\mathtt{val}}
\newcommand*{\Letm}{\mathtt{mod}}

\newcommand*{\ValRel}[1]{\mathcal{V}\sembracket{#1}}
\newcommand*{\ExprRel}[1]{\mathcal{E}\sembracket{#1}}
\newcommand*{\CtxRel}[1]{\mathcal{C}\sembracket{#1}}
\newcommand*{\ModRel}[1]{\mathcal{M}\sembracket{#1}}
\newcommand*{\TyEnv}{\text{TyEnv}}
\newcommand*{\TyVar}{\text{TyVar}}
\newcommand*{\Type}{\text{Type}}
\newcommand*{\Subst}{\text{Subst}}
\newcommand*{\external}{\Gamma_{\text{ext}}}

\newcommand*{\InitE}{\textsf{Init}}
\newcommand*{\ReadE}{\textsf{Read}}
\newcommand*{\CallE}{\textsf{Call}}
\newcommand*{\resolve}{\downarrow_{\ctx_0}}


% ====================
% Lengths
% ====================

% If you have N columns, choose \sepwidth and \colwidth such that
% (N+1)*\sepwidth + N*\colwidth = \paperwidth
\newlength{\sepwidth}
\newlength{\colwidth}
\setlength{\sepwidth}{0.025\paperwidth}
\setlength{\colwidth}{0.4625\paperwidth}

\newcommand{\separatorcolumn}{\begin{column}{\sepwidth}\end{column}}

% ====================
% Title
% ====================

\title{프로그램 따로분석의 이론적 기틀: 재귀적 모듈 지원하기}

\author{이준협 \and 이광근}

\institute[shortinst]{서울대학교 프로그래밍 연구실 (ROPAS)}

% ====================
% Footer (optional)
% ====================

%\footercontent{
%  \hfill
%  2023년 SIGPL 여름학교 포스터 발표
%  \hfill
%}
% (can be left out to remove footer)

% ====================
% Logo (optional)
% ====================

% use this to include logos on the left and/or right side of the header:
\logoright{\includegraphics[height=5cm]{ropas-symbol.png}}
\logoleft{\includegraphics[height=5cm]{snu-symbol.png}}

% ====================
% Body
% ====================

\begin{document}

\begin{frame}[t]
  \begin{columns}[t]
    \separatorcolumn

    \begin{column}{\colwidth}
      \begin{block}{문제 소개}
        \heading{무엇을 하고 싶은가?}
        \begin{itemize}
          \item 프로그램 외부의 값을 몰라도 분석하고 싶다.\\
                \textbf{ex. }외부 라이브러리 함수를 부르는 경우
          \item 이름이 ``복잡한 방식''으로 알려질 때에도 분석하고 싶다.\\
                \textbf{ex. }First-class modules, Functors, Recursive modules
        \end{itemize}
        \heading{어떻게 할 것인가?}
        \begin{itemize}
          \item ``Algebraic Effects'': 환경을 받는 \textbf{함수}를 \textbf{자료구조}로 나타냄.
          \item 환경과의 상호작용을 요약하고, 나중에 환경이 들어오면 풀자!
        \end{itemize}
      \end{block}

      \begin{block}{모듈이 있는 언어의 정의}
        \heading{겉모습 (Untyped $\lambda$+Modules)}
        \begin{figure}[h!]
          \large
          \begin{tabular}{rrcll}
            Identifiers & $x$ & $\in$         & $\ExprVar$                                                      \\
            Expression  & $e$ & $\rightarrow$ & $x$ $\vbar$ $\lambda x.e$ $\vbar$ $e$ $e$ & $\lambda$-calculus  \\
                        &     & $\vbar$       & $\link{e}{e}$                             & linked expression   \\
                        &     & $\vbar$       & $\varepsilon$                             & empty module        \\
                        &     & $\vbar$       & $x=e$ ; $e$                               & (recursive) binding
          \end{tabular}
        \end{figure}
        \heading{속내용}
        \begin{figure}[h!]
          \centering
          \large
          \begin{tabular}{rrcll}
            Environment & $\ctx$ & $\in$         & $\Ctx$                                                                        \\
            Location    & $\ell$ & $\in$         & $\Loc$                                                                        \\
            Value       & $v$    & $\in$         & $\Value \triangleq\Ctx+\ExprVar\times\Expr\times\Ctx$                         \\
            Weak Value  & $w$    & $\in$         & $\Walue\triangleq\Value+\Loc\times\Value$                                     \\
            Environment & $\ctx$ & $\rightarrow$ & $\bullet$                                             & empty stack           \\
                        &        & $\vbar$       & $(x,w)\cons \ctx$                                     & weak value binding    \\
                        &        & $\vbar$       & $(x,\ell)\cons\ctx$                                   & free location binding \\
            Value       & $v$    & $\rightarrow$ & $\ctx$                                                & exported environment  \\
                        &        & $\vbar$       & $\langle \lambda x.e, \ctx \rangle$                   & closure               \\
            Weak Value  & $w$    & $\rightarrow$ & $v$                                                   & value                 \\
                        &        & $\vbar$       & $\mu\ell.v$                                           & recursive value
          \end{tabular}
        \end{figure}
        \heading{실행의미}
        \begin{figure}[h!]
          \centering\large
          \begin{flushright}
            \fbox{$(e,\ctx)\Downarrow v$}
          \end{flushright}
          \vspace{0pt} % -0.75em}
          \[
            \begin{InfRule}{Id}
              \hypo{\ctx(x)=v}
              \infer1{
                (x, \ctx)
                \Downarrow
                v
              }
            \end{InfRule}\quad
            \fbox{
              \begin{InfRule}{RecId}
                \hypo{\ctx(x)=\mu\ell.v}
                \infer1{
                  (x, \ctx)
                  \Downarrow
                  v[\mu\ell.v/\ell]
                }
              \end{InfRule}}\quad
            \begin{InfRule}{Fn}
              \infer0{
                (\lambda x.e, \ctx)
                \Downarrow
                \langle\lambda x.e, \ctx\rangle
              }
            \end{InfRule}
          \]

          \[
            \begin{InfRule}{App}
              \hypo{
                (e_{1}, \ctx)
                \Downarrow
                \langle\lambda x.e, \ctx_1\rangle
              }
              \hypo{
                (e_{2}, \ctx)
                \Downarrow
                v_2
              }
              \hypo{
                (e, (x, v_2)\cons \ctx_1)
                \Downarrow
                v
              }
              \infer3{
                (e_{1}\:e_{2}, \ctx)
                \Downarrow
                v
              }
            \end{InfRule}
          \]

          \[
            \fbox{
              \begin{InfRule}{Link}
                \hypo{
                  (e_{1}, \ctx)
                  \Downarrow
                  \ctx_1
                }
                \hypo{
                  (e_{2}, \ctx_1)
                  \Downarrow
                  v
                }
                \infer2{
                  (\link{e_{1}}{e_{2}}, \ctx)
                  \Downarrow
                  v
                }
              \end{InfRule}}\quad
            \begin{InfRule}{Empty}
              \infer0{
                (\varepsilon, \ctx)
                \Downarrow
                \bullet
              }
            \end{InfRule}
          \]

          \[
            \fbox{
              \begin{InfRule}{RecBind}
                \hypo{
                  \ell\not\in\FLoc(\ctx)
                }
                \hypo{
                  (e_{1}, (x,\ell)\cons\ctx)
                  \Downarrow
                  v_1
                }
                \hypo{
                  (e_{2}, (x, \mu\ell.v_1)\cons \ctx)
                  \Downarrow
                  \ctx_2
                }
                \infer3{
                  (x=e_1; e_2, \ctx)
                  \Downarrow
                  (x,\mu\ell.v_1)\cons\ctx_2
                }
              \end{InfRule}}
          \]
        \end{figure}
      \end{block}

    \end{column}

    \separatorcolumn

    \begin{column}{\colwidth}
      \begin{block}{상호작용 기록하고 나중에 풀기}
        \heading{상호작용 기록}
        \begin{figure}[h!]
          \centering
          \large
          \begin{tabular}{rrcll}
            Environment & $\ctx$ & $\rightarrow$ & $\cdots$                            \\
                        &        & $\vbar$       & $[E]$         & answer to an event  \\
            Value       & $v$    & $\rightarrow$ & $\cdots$                            \\
                        &        & $\vbar$       & $E$           & answer to an event  \\
            Event       & $E$    & $\rightarrow$ & $\InitE$      & initial environment \\
                        &        & $\vbar$       & $\ReadE(E,x)$ & read event          \\
                        &        & $\vbar$       & $\CallE(E,v)$ & call event
          \end{tabular}

          \[
            [E](x)\triangleq\ReadE(E,x)\qquad
            \begin{InfRule}{AppEvent}
              \hypo{
                (e_{1}, \ctx)
                \Downarrow
                E
              }
              \hypo{
                (e_{2}, \ctx)
                \Downarrow
                v
              }
              \infer2{
                (e_{1}\:e_{2}, \ctx)
                \Downarrow
                \CallE(E,v)
              }
            \end{InfRule}
          \]
        \end{figure}

        \heading{상호작용 풀기 (일부 규칙들)}
        \begin{figure}[h!]
          \centering
          \large
          \begin{flushright}
            \fbox{$w\resolve W$}
          \end{flushright}
          \vspace{0pt} % -0.75em}
          \[
            \begin{InfRule}{R-Init}
              \infer0{
                \InitE
                \downarrow
                \ctx_0
              }
            \end{InfRule}\quad
            \begin{InfRule}{R-Read}
              \hypo{E\downarrow\Sigma}
              \hypo{\Sigma(x)=V}
              \infer2{
                \ReadE(E,x)
                \downarrow
                V
              }
            \end{InfRule}\quad
            \begin{InfRule}{R-ReadRec}
              \hypo{E\downarrow\Sigma}
              \hypo{\Sigma(x)=\mu\ell.V}
              \infer2{
                \ReadE(E,x)
                \downarrow
                V[\mu\ell.V/\ell]
              }
            \end{InfRule}
          \]

          \[
            \begin{InfRule}{R-CallV}
              \hypo{E\downarrow\langle\lambda x.e,\Sigma\rangle}
              \hypo{v\downarrow V}
              \hypo{(e,(x,V)\cons\Sigma)\semarrow V'}
              \infer3{
                \CallE(E,v)
                \downarrow
                V'
              }
            \end{InfRule}\quad
            \begin{InfRule}{R-CallE}
              \hypo{E\downarrow E'}
              \hypo{v\downarrow V}
              \infer2{
                \CallE(E,v)
                \downarrow
                \CallE(E',V)
              }
            \end{InfRule}
          \]
        \end{figure}

        \heading{Conjecture (상호작용 풀기는 안전하다)}
        \[\ctx\resolve\Sigma\text{ and }(e,\ctx)\semarrow v\text{ and }(e,\Sigma)\semarrow V\Rightarrow v\resolve V\]

        \heading{예시}
        {\ttfamily
          \begin{tabular}{rl}
            (\textcolor{dkviolet}{Top}=          & (\textcolor{dkgreen}{Tree}=(\textcolor{blue}{max}=$\lambda$x.f (Top$\synlink$Forest$\synlink$max) x; \textcolor{blue}{$\varepsilon$}); \\
                                                 & \textcolor{dkgreen}{Forest}=(\textcolor{red}{max}=$\lambda$x.(Top$\synlink$Tree$\synlink$max) x; \textcolor{red}{$\varepsilon$});      \\
                                                 & \textcolor{dkgreen}{$\varepsilon$});                                                                                                   \\
            \textcolor{dkviolet}{$\varepsilon$}) & $\synlink$((\textcolor{dkviolet}{Top}$\synlink$\textcolor{dkgreen}{Tree}$\synlink$\textcolor{blue}{max}) id)
          \end{tabular}
        }

        $\semarrow\CallE(\CallE(\ReadE(\InitE,\text{``f''}),\langle\lambda\texttt{x.Top}\synlink\texttt{Tree}\synlink\texttt{max x},\dots\rangle),\texttt{id})$
      \end{block}
      \begin{block}{요약하기}
        \heading{메모리의 등장}
        \begin{itemize}
          \item 메모리가 있으면 요약이 편해진다.
          \item $\Ctx\triangleq\fin{\ExprVar}{\Loc}$, $\Mem\triangleq\fin{\Loc}{\Value}$이면, $\Loc$만 요약.
          \item 그렇다면 실행의미가 동일하다는 것부터 보여야 한다.
        \end{itemize}
        \begin{align*}
          \ctx\sim\ctx_0,\mem_0\land(e,\ctx)\semarrow v                    & \Rightarrow\exists v_1,\mem_1\sim v:(e,\ctx_0,\mem_0)\semarrow(v_1,\mem_1) \\
          \ctx\sim\ctx_0,\mem_0\land(e,\ctx_0,\mem_0)\semarrow(v_1,\mem_1) & \Rightarrow\exists v\sim v_1,\mem_1:(e,\ctx)\semarrow v
        \end{align*}
        \heading{$\semarrow$와 $\resolve$의 요약 (진행 중)}
        \begin{itemize}
          \item 첫번째: $\downarrow$를 메모리가 있을 때 ``안전''하게 정의하기
                \[\text{``안전'':}w\sim v,\mem\land w\downarrow W\Rightarrow\exists V,M:W\sim V,M\land v,\mem\downarrow V,M\]
          \item 두번째: 메모리의 요약에 대해 ``안전''하게 $\downarrow$, $\Downarrow$를 요약
        \end{itemize}
      \end{block}
    \end{column}

    \separatorcolumn
  \end{columns}
\end{frame}

\end{document}
